\section{Objective}
\begin{itemize}
	\item To design, construct, and evaluate a full-wave bridge rectifier circuit on a breadboard using silicon diodes.
	\item To examine the cycle-by-cycle conduction path during positive and negative input AC half-cycles.
	\item To record input and output voltage waveforms using a Digital Storage Oscilloscope (DSO).
	\item To evaluate key performance parameters, including average DC output voltage ($V_{DC}$), RMS output voltage ($V_{rms}$), ripple factor ($r$), and rectification efficiency ($\eta$).
	\item To assess output smoothing achieved via parallel capacitive filters under varying load conditions.
\end{itemize}

\section{Introduction}
Rectification convert bidirectional Alternating Current (AC) into unidirectional Direct Current (DC). Among standard topologies, full-wave bridge rectifiers offer superior utilization of AC supply cycles without requiring a center-tapped transformer. Arranging four diodes in a bridge array directs current through the load resistor in a single physical direction across both input half-cycles. The output waveform is subsequently smoothed using parallel capacitive filtering to deliver stable DC levels appropriate for electronic circuitry.

\section{Theory}

\subsection{Conduction Analysis}
Consider the AC voltage delivered by the step-down transformer:
\[
v_s(t) = V_m \sin(\omega t),
\]
where $V_m = V_{rms}\sqrt{2}$ is peak secondary voltage and $\omega = 2\pi f$ is the supply angular frequency.
\begin{itemize}
	\item \textbf{Positive Half-Cycle ($0 < \omega t < \pi$):} The upper secondary terminal operates at higher potential relative to the lower terminal. Diodes $D_1$ and $D_4$ are forward-biased, while $D_2$ and $D_3$ remain reverse-biased. Current travels through $D_1$, downward through load $R_L$, and returns via $D_4$.
	\item \textbf{Negative Half-Cycle ($\pi < \omega t < 2\pi$):} Supply polarity reverses. Diodes $D_2$ and $D_3$ transition to forward conduction, whereas $D_1$ and $D_4$ become reverse-biased. Current passes through $R_L$ in the identical physical direction (top-to-bottom) via $D_2$ and $D_3$.
\end{itemize}
Because current flows through the load in both half-cycles, the output ripple frequency is double the input frequency ($f_{\text{out}} = 2 f_{\text{in}}$).

\subsection{Mathematical Formulation}
Accounting for the forward voltage drop ($V_D \approx \qty{0.7}{\volt}$) across two conducting diodes in each half-cycle, peak load voltage is:
\[
V_{mo} = V_m - 2V_D.
\]
The average DC output voltage and RMS output voltage are calculated as:
\[
V_{DC} = \frac{2V_{mo}}{\pi}, \qquad V_{rms} = \frac{V_{mo}}{\sqrt{2}}.
\]
The ripple factor ($r$), representing relative AC content in the rectified output, is expressed by:
\[
r = \sqrt{\left(\frac{V_{rms}}{V_{DC}}\right)^2 - 1} = \sqrt{\left(\frac{\pi}{2\sqrt{2}}\right)^2-1} \approx 0.482.
\]
Theoretical maximum rectification efficiency and peak inverse voltage (PIV) rating per diode are:
\[
\eta_{\max} = \frac{8}{\pi^2}\approx 81.2\%, \qquad \text{PIV} = V_m.
\]

\subsection{Parallel Capacitive Filter}
Connecting a filter capacitor $C$ in parallel with load $R_L$ smooths the output profile by charging rapidly near voltage peaks and discharging through the load between cycles. Peak-to-peak ripple voltage ($V_{r(p\text{-}p)}$) and overall ripple factor ($r$) are modeled as:
\[
V_{r(p\text{-}p)} \approx \frac{I_{DC}}{2fC} = \frac{V_{DC}}{2fR_LC}, \qquad r \approx \frac{1}{4\sqrt{3}\,fR_LC}.
\]
Increasing either capacitance $C$ or load resistance $R_L$ minimizes output voltage ripple.

\section{Apparatus and Components}
\begin{table}[h!]
	\centering
	\caption{List of Equipment and Components Used}
	\vspace{2mm}
	\begin{tabular}{llc}
		\toprule
		\rowcolor{blue!10}
		\textbf{Item Description} & \textbf{Specification / Rating} & \textbf{Quantity} \\
		\midrule
		Step-Down Transformer & \qty{230}{\volt}/\qty{12}{\volt} RMS, \qty{50}{\hertz} & 1 \\
		Silicon Rectifier Diodes & 1N4007 & 4 \\
		Load Resistor $R_L$ & \qty{1}{\kilo\ohm}, \qty{2.2}{\kilo\ohm} & 2 \\
		Electrolytic Filter Capacitor $C$ & \qty{47}{\micro\farad}, \qty{100}{\micro\farad} & 2 \\
		Digital Multimeter (DMM) & DC/AC Volts & 1 \\
		Digital Storage Oscilloscope (DSO) & Dual-Channel with 10$\times$ Probes & 1 \\
		Solderless Breadboard & Standard Tie-Point & 1 \\
		Connecting Wires & Jumper Wires & As required \\
		\bottomrule
	\end{tabular}
	\label{tab:apparatus6}
\end{table}

\section{Circuit Diagram}
\begin{figure}[h!]
	\centering
	\begin{circuitikz}[american, scale=0.85]
		\draw (-2,0) to[sV, l=$v_s(t)$] (-2,4);
		\node at (-2,-0.6) [below] {Transformer Secondary};
		\draw (-2,4) -- (0,4) -- (0,2);
		\draw (-2,0) -- (-2,-3) -- (6,-3) -- (6,2);
		\coordinate (L) at (0,2);
		\coordinate (R) at (6,2);
		\coordinate (T) at (3,5);
		\coordinate (Bt) at (3,-1);
		\draw (L) to[D, l=$D_1$] (T);
		\draw (R) to[D, l=$D_2$] (T);
		\draw (Bt) to[D, l=$D_3$] (L);
		\draw (Bt) to[D, l=$D_4$] (R);
		\draw (T) to[R, l=$R_L$] (Bt);
		\draw (T) -- (4.5,5) to[C, l=$C$] (4.5,-1) -- (Bt);
		\draw (Bt) node[ground]{};
		\node at (0,2.4) [above left] {AC$_1$};
		\node at (6,2.4) [above right] {AC$_2$};
		\node at (3,5.4) [above] {$+$ (DC Output)};
	\end{circuitikz}
	\caption{Full-wave bridge rectifier circuit incorporating transformer feed, four-diode bridge array, load resistor $R_L$, and optional parallel filter capacitor $C$.}
	\label{fig:bridgerectifier}
\end{figure}

\section{Experimental Procedure}
\begin{enumerate}
	\item Construct the bridge configuration using four 1N4007 diodes on the breadboard, observing correct terminal polarity.
	\item Connect the \qty{12}{\volt} secondary terminals of the transformer to the AC input terminals ($\text{AC}_1$, $\text{AC}_2$) of the bridge network.
	\item Connect $R_L=\qty{1}{\kilo\ohm}$ across the DC output without capacitive filtering. Observe AC input and rectified DC output signals on Channel 1 and Channel 2 of the DSO respectively. Measure $V_{DC}$ and $V_{rms}$ using the DMM.
	\item Connect filter capacitor $C=\qty{47}{\micro\farad}$ across $R_L$. Record output $V_{DC}$ and peak-to-peak ripple voltage $V_{r(p\text{-}p)}$ on the DSO.
	\item Repeat filtered measurements for $C=\qty{100}{\micro\farad}$ and load setting $R_L=\qty{2.2}{\kilo\ohm}$ to evaluate ripple suppression behavior.
\end{enumerate}

\section{Theoretical Calculation}
Using supply voltage $V_{rms}=\qty{12}{\volt}$, frequency $f=\qty{50}{\hertz}$, diode drop $V_D=\qty{0.7}{\volt}$, and load $R_L=\qty{1}{\kilo\ohm}$:
\[
V_m = V_{rms}\sqrt{2} = 12 \times 1.4142 = \qty{16.97}{\volt},
\]
\[
V_{mo} = V_m - 2V_D = 16.97 - 1.4 = \qty{15.57}{\volt},
\]
\[
V_{DC} = \frac{2V_{mo}}{\pi} = \frac{2 \times 15.57}{3.1416} = \qty{9.91}{\volt}, \qquad
I_{DC} = \frac{V_{DC}}{R_L} = \frac{9.91}{1000} = \qty{9.91}{\milli\ampere},
\]
\[
V_{rms(\text{out})} = \frac{V_{mo}}{\sqrt{2}} = \frac{15.57}{1.4142} = \qty{11.01}{\volt},
\]
\[
\eta = \frac{V_{DC}^2 / R_L}{V_{rms(\text{out})}^2 / R_L} \times 100\% = \frac{9.91^2}{11.01^2} \times 100\% \approx 81.0\%, \qquad \text{PIV} = V_m = \qty{16.97}{\volt}.
\]
With filter capacitor $C = \qty{47}{\micro\farad}$ connected:
\[
V_{r(p\text{-}p)} \approx \frac{V_{DC}}{2fR_LC} = \frac{9.91}{2 \times 50 \times 1000 \times 47 \times 10^{-6}} \approx \qty{2.11}{\volt}.
\]

\section{Data Collection}

\begin{table}[h!]
	\centering
	\caption{Unfiltered Rectifier Performance Across Load Settings}
	\vspace{2mm}
	\resizebox{\textwidth}{!}{%
		\begin{tabular}{ccccccc}
			\toprule
			\rowcolor{blue!10}
			\textbf{Load $R_L$} & \textbf{Peak Input $V_m$} & \textbf{Peak Output $V_{mo}$} & \textbf{Measured $V_{DC}$} & \textbf{Measured $V_{rms}$} & \textbf{Ripple Factor $r$} & \textbf{Efficiency $\eta$} \\
			\midrule
			\qty{1}{\kilo\ohm} & \textit{TBD} & \textit{TBD} & \textit{TBD} & \textit{TBD} & \textit{TBD} & \textit{TBD} \\
			\qty{2.2}{\kilo\ohm} & \textit{TBD} & \textit{TBD} & \textit{TBD} & \textit{TBD} & \textit{TBD} & \textit{TBD} \\
			\bottomrule
		\end{tabular}%
	}
	\label{tab:unfiltereddata6}
\end{table}

\begin{table}[h!]
	\centering
	\caption{Filtered Rectifier Performance: Smoothing Effect of $C$}
	\vspace{2mm}
	\begin{tabular}{cccc}
		\toprule
		\rowcolor{blue!10}
		\textbf{Load $R_L$} & \textbf{Filter $C$} & \textbf{Measured $V_{DC}$} & \textbf{Ripple $V_{r(p\text{-}p)}$} \\
		\midrule
		\qty{1}{\kilo\ohm} & \qty{47}{\micro\farad} & \textit{TBD} & \textit{TBD} \\
		\qty{1}{\kilo\ohm} & \qty{100}{\micro\farad} & \textit{TBD} & \textit{TBD} \\
		\qty{2.2}{\kilo\ohm} & \qty{47}{\micro\farad} & \textit{TBD} & \textit{TBD} \\
		\bottomrule
	\end{tabular}
	\label{tab:filtereddata6}
\end{table}

\begin{tcolorbox}[colback=gray!5, colframe=primaryColor, arc=2mm, title=\textbf{Note on Experimental Data}]
	Laboratory measurements will be populated upon final validation of the experimental data sheet. All formatting parameters and analytical frameworks are fully prepared to evaluate collected parameters.
\end{tcolorbox}

\section{Result Analysis}
Experimental readings will be compared against theoretical formulations using percentage error:
\[
\text{\% Error} = \left|\frac{\text{Theoretical} - \text{Measured}}{\text{Theoretical}}\right| \times 100\%.
\]
Output ripple frequency will be evaluated against $f_{\text{out}} = 2f_{\text{in}} = \qty{100}{\hertz}$. The reduction in peak-to-peak ripple voltage resulting from higher capacitance values ($C = \qty{100}{\micro\farad}$) will be verified against theoretical models.

\section{Discussion and Sources of Error}
Deviations between experimental and theoretical parameters stem from the following technical factors:
\begin{enumerate}
	\item \textbf{Diode Forward Drop Variations:} Actual forward voltage drop across 1N4007 diodes under load may vary from the nominal \qty{0.7}{\volt} model.
	\item \textbf{Transformer Regulation Loss:} Internal transformer winding resistance causes output voltage droop under load conditions.
	\item \textbf{Capacitor Tolerances and ESR:} Electrolytic capacitor tolerance ranges and non-zero Equivalent Series Resistance (ESR) alter peak ripple attenuation.
	\item \textbf{Measurement Equipment Loading:} DSO probe input impedance and minor AC line frequency variations introduce small experimental errors.
\end{enumerate}

\section{Conclusion}
The design and operational mechanics of the full-wave bridge rectifier were systematically evaluated. The circuit effectively doubles ripple frequency relative to input AC supply and increases DC conversion efficiency. The integration of parallel capacitive filtering demonstrates substantial attenuation of AC ripple components, delivering a stabilized DC voltage profile.

\section{References}
\begin{tcolorbox}[colback=gray!5, colframe=gray!40, arc=2mm]
	\begin{enumerate}
		\item Boylestad, R. L., \& Nashelsky, L. (2013). \textit{Electronic Devices and Circuit Theory} (11th ed.). Pearson.
		\item Alexander, C. K., \& Sadiku, M. N. O. (2016). \textit{Fundamentals of Electric Circuits} (6th ed.). McGraw-Hill Education.
		\item Sedra, A. S., \& Smith, K. C. (2014). \textit{Microelectronic Circuits} (7th ed.). Oxford University Press.
	\end{enumerate}
\end{tcolorbox}